% =====================================================================
% Section 00 — Front matter (Abstract)
% CONTENT PLAN: one tight abstract. State the question (what does scale buy),
%   the 3 measured findings (operand saturates; load retires in order;
%   compliance per-write rate flat), the cost payoff (clearly an estimate),
%   and the conclusion (scale buys capability not guarantee -> cost knee).
%   No internal tags, no projections stated as results.
% SUBSECTION TITLES: n/a (abstract has none).
% =====================================================================
\begin{abstract}
Tool-use agents map a natural-language request to a sequence of tool calls under a domain policy, and are usually improved by scaling the base model. We ask instead, function by function and across scale, what scaling actually buys: which sub-capabilities improve with size, which remain limited at every scale, which are better supplied by deterministic scaffolding or by per-domain authored relations (which we call A2), and how these should be assembled to minimize total deployment cost. Treating four benchmarks ($\tau^2$-bench, SOPBench, TaskBench, and Synth) as measurement instruments, we decompose tool-use failure into \emph{operand execution}, \emph{orchestration load}, and \emph{compliance}, and chart each across model scale (7B--32B measured; 72B and above projected) and across levers (base model, capability scaffold, guarantee scaffold, A2, and learning). We report three findings. First, operand execution saturates early: given an explicit specification, a 32B base model selects the correct action on 88 of 88 cases, so this is not a learnable gap at that scale. Second, the load dimensions that predict failure retire with scale in a definite order---interference and state-multiplicity bind at 7B and 14B but disappear by 32B, while context length and conditional depth persist---and a controlled probe overturns an interference signal that appears only in observational correlations. Third, compliance is scale-invariant: task success rises with scale (0.19 at 7B to 0.55 at 32B), yet the per-write-opportunity rate of policy-semantic violations does not fall ($\approx$0.07--0.10 across all measured scales, with overlapping 95\% confidence intervals), even as authorization-type violations collapse; a deterministic policy gate eliminates the residual at every scale. We interpret deterministic scaffolding as a load-reduction operator and fold the map into a total-lifecycle cost objective. On $\tau^2$-retail, a 32B int8 base with a deterministic gate reaches near-equal compliance to a frontier API while running on-premises; a rough token-based estimate puts its per-request cost materially lower, though the exact multiple is deployment-specific and not measured, and an on-premises-plus-escalation fleet could in principle reach equal-or-better compliance at lower blended cost. We conclude that scale buys capability but not guarantee: even a frontier model requires a deterministic guarantee layer, so the cost-optimal design is not the smallest model but the total-cost knee.
\end{abstract}
